\documentclass[12pt]{article}
\usepackage[utf8]{inputenc}
\usepackage[T1]{fontenc}
\usepackage{hyperref}
\title{Puma Language Resources for K-12 Curriculum Development}
\author{Compiled by Claude Sonnet 5.5 (claude-sonnet-5-5) for Andr\'as Kornai}
\date{}
\begin{document}
\maketitle

\section{Introduction}
Puma is a Kiranti (Sino-Tibetan) language of eastern Nepal. The 2001 census counted about 4,310 speakers and the 2011 census about 6,700. It has been documented by the Chintang and Puma Documentation Project, which reports that children still learn the language and that efforts are under way to bring it into school education. Puma is distinct from Chamling, another Kiranti language.

\section{Primary Educational Resources}
No school textbooks or primers were found in the searches. The documentation project site states an aim of supporting school use.

\section{Linguistic and Reference Materials}
\begin{itemize}
\item Bickel, Gaenszle et al. (2009), a Puma--Nepali dictionary and grammar published in Lalitpur by Kir\=at Pum\=a R\=a\=i Tupkh\=ab\=a\.nkh\=al\=a. The record in the University of Zurich repository (its DOI link appears unrelated and should be checked). \url{https://www.zora.uzh.ch/entities/publication/03328fc1-de0b-4965-ac3b-4766a9fa8b00}
\item Narayan Prasad Sharma (2014), ``Morphosyntax of Puma'', SOAS PhD thesis, mentioned in search summaries; a direct link was not seen.
\item Wikipedia entry on the Puma language. \url{https://wikipedia.com/wiki/ISO_639:pum}
\item Endangered Languages Project bibliography for Puma. \url{https://endangeredlanguages.com/lang/1125/bibliography}
\end{itemize}

\section{Religious and Community Resources}
\begin{itemize}
\item Kirat Rai Chamling Language Research and Investigation (a related but different Kiranti language, listed for regional context only). \url{https://kiratrai.org/?p=1777}
\end{itemize}

\section{Digital and Audio Resources}
\begin{itemize}
\item Chintang and Puma Documentation Project (Nepali-language site). \url{https://www.uzh.ch/cpdp/left%20home%20frame.html}
\item DoBeS poster on Chintang and Puma documentation. \url{https://dobes.mpi.nl/posters/print/eurasia-chintang-print.pdf}
\end{itemize}

\section{Curriculum Development Context}
Puma has a published dictionary and grammar and an extensive documentation project, which makes it better resourced than many comparable languages. The next step would be to contact the project and the community association about the corpus and any school materials.
\end{document}